\documentclass[11pt,letterpaper]{article}
\usepackage[top=0.9in,bottom=0.9in,left=1in,right=1in]{geometry}
\usepackage{amsmath,amssymb}
\usepackage{xcolor}
\usepackage{mdframed}
\usepackage{enumitem}
\usepackage{booktabs}
\usepackage{fancyhdr}

% ── Color scheme matching the cheat sheet ──
\definecolor{purple}{HTML}{534AB7}\definecolor{purplebg}{HTML}{EEEDFE}
\definecolor{teal}{HTML}{0F6E56}\definecolor{tealbg}{HTML}{E1F5EE}
\definecolor{coral}{HTML}{993C1D}\definecolor{coralbg}{HTML}{FAECE7}
\definecolor{green}{HTML}{3B6D11}\definecolor{greenbg}{HTML}{EAF3DE}
\definecolor{blue}{HTML}{0C447C}\definecolor{bluebg}{HTML}{E6F1FB}
\definecolor{amber}{HTML}{854F0B}\definecolor{amberbg}{HTML}{FAEEDA}

% Cheat sheet cross-reference tag — shows which page and section to look up
\newcommand{\cs}[1]{\textcolor{teal}{\footnotesize[CS\,#1]}}

% Highlighted note box
\newmdenv[linewidth=1pt,linecolor=coral,backgroundcolor=coralbg!40,
  innerleftmargin=6pt,innerrightmargin=6pt,
  innertopmargin=4pt,innerbottommargin=4pt,
  skipabove=4pt,skipbelow=4pt]{notebox}

\pagestyle{fancy}\fancyhf{}
\renewcommand{\headrulewidth}{0.4pt}
\fancyhead[L]{\textbf{ECE332 Fall 2025 --- Exam 1 Answer Key (Corrected)}}
\fancyhead[R]{Bao Nguyen}
\fancyfoot[C]{\thepage}

\setlength{\parindent}{0pt}
\setlength{\parskip}{4pt}

\newcommand{\step}[1]{\medskip\noindent\textbf{Step #1.}\enspace}

\begin{document}

\begin{center}
  {\Large\textbf{ECE 332 Fall 2025 --- Exam 1 Answer Key}}\\[4pt]
  {\small \textcolor{teal}{[CS\,p.1: SECTION]} = page 1 of cheat sheet;\enspace
   \textcolor{teal}{[CS\,p.2: SECTION]} = page 2.}
\end{center}

\begin{notebox}
\textbf{\textcolor{coral}{Two corrected errors from prior answer keys:}}\\
\textbf{(1) Q1(c):} The exam writes ``$\omega=2\,\text{MHz}$'' but MHz is a \emph{frequency} unit.
Must convert first: $\omega = 2\pi f = 2\pi\times2\times10^6 = 4\pi\times10^6\,\text{rad/s}$.
Both prior keys skipped this step and used $2\times10^6$ directly.\\
\textbf{(2) Q6:} The exam gives $f=\tfrac{2}{\pi}\,\text{GHz}$, not $\tfrac{2}{5}\,\text{GHz}$.
$\omega = 2\pi f = 2\pi\cdot\tfrac{2}{\pi}\times10^9 = 4\times10^9\,\text{rad/s}$.
\end{notebox}

\vspace{6pt}

%──────────────────────────────────────────────────────────────────────
\colorbox{purplebg}{\textbf{\color{purple}\large Question 1 --- Triangular Loop (20 pts)}}
%──────────────────────────────────────────────────────────────────────

\vspace{4pt}
\textbf{Setup.} Equilateral loop in $y$-$z$ plane, $N$ turns, side $a$.
$\mathbf{B}(t)=B_0(-2\hat{x}+3\hat{y})\sin(\omega t)$.

\medskip

\textbf{(a) Magnetic flux through one turn.} \hfill\textit{(6 pts)}

\cs{p.1: SURFACE ELEMENT $d\mathbf{s}$}: loop in $y$-$z$ plane $\Rightarrow$ $d\mathbf{s}=\hat{x}\,da$.\\
\cs{p.1: DOT PRODUCTS}: $\hat{x}\cdot\hat{x}=1$, $\hat{y}\cdot\hat{x}=0$
$\Rightarrow$ only the $-2\hat{x}$ component contributes.\\
\cs{p.1: ROTATING LOOP area table}: equilateral triangle with side $a$: $A=\tfrac{\sqrt{3}}{4}a^2$.

\[
\Phi = \int_S \mathbf{B}\cdot d\mathbf{s}
= B_0(-2\underbrace{\hat{x}\cdot\hat{x}}_{1}+3\underbrace{\hat{y}\cdot\hat{x}}_{0})\sin(\omega t)\cdot A
\]
\[
\boxed{\Phi = -\tfrac{\sqrt{3}}{2}\,B_0\,a^2\sin(\omega t)\;\text{Wb}}
\]

\medskip

\textbf{(b) EMF induced in the loop.} \hfill\textit{(6 pts)}

\cs{p.1: EMF \& FARADAY'S LAW}: $V_\text{emf}=-N\dfrac{d\Phi}{dt}$.\\
\cs{p.1: $\partial B/\partial t$ TABLE}: $\dfrac{d}{dt}[\sin(\omega t)]=\omega\cos(\omega t)$.

\[
V_\text{emf} = -N\,\frac{d}{dt}\!\left[-\tfrac{\sqrt{3}}{2}B_0 a^2\sin(\omega t)\right]
\]
\[
\boxed{V_\text{emf} = \frac{\sqrt{3}}{2}\,N B_0\,\omega\,a^2\cos(\omega t)\;\text{V}}
\]

\medskip

\textbf{(c) Current through 1\,k$\Omega$ at given values.} \hfill\textit{(8 pts)}

\begin{notebox}
\textbf{Unit conversion required.}
The exam states ``$\omega=2\,\text{MHz}$'', but MHz is a frequency unit (cycles/s).
\cs{p.1: CIRCUIT \& WAVE strip}: $\omega = 2\pi f$, so
\[
\omega = 2\pi \times 2\times10^6 = 4\pi\times10^6\;\text{rad/s}
\]
\cs{p.1: SI PREFIXES}: $\text{MHz}=10^6$.
\end{notebox}

Given: $N=10$, $a=0.08\,\text{m}$, $B_0=0.02\,\text{T}$,
$\omega=4\pi\times10^6\,\text{rad/s}$, $\omega t=\pi/4$, $R=1\,\text{k}\Omega$.

\step{1} Combine trig factors.\\
\cs{p.1: UNIT CIRCLE}: $\cos(\pi/4)=\tfrac{\sqrt{2}}{2}$.
\[
\frac{\sqrt{3}}{2}\cdot\frac{\sqrt{2}}{2} = \frac{\sqrt{6}}{4}
\]

\step{2} Compute $|V_\text{emf}|$ numerically.
\begin{align*}
|V_\text{emf}| &= \frac{\sqrt{6}}{4}\cdot N\cdot B_0\cdot\omega\cdot a^2\\
&= \frac{\sqrt{6}}{4}\cdot10\cdot0.02\cdot(4\pi\times10^6)\cdot(0.08)^2
\end{align*}
Bracket step by step:
\[
(0.08)^2 = 6.4\times10^{-3},\quad
0.02\times6.4\times10^{-3} = 1.28\times10^{-3}
\]
\[
1.28\times10^{-3}\times4\pi\times10^6 = 5.12\pi\times10^3,\quad
10\times5.12\pi\times10^3 = 5120\pi\times10^1
\]
Wait, step more carefully:
\[
10\times 1.28\times10^{-3}\times4\pi\times10^6
= 10\times5.12\pi\times10^3
= 51200\pi
\]
\[
|V_\text{emf}| = \frac{\sqrt{6}}{4}\times 51200\pi = 12800\pi\sqrt{6}
\approx 12800\times3.1416\times2.449
\approx 98{,}490\;\text{V}
\]

Hmm — let me redo carefully:
\[
10\times0.02 = 0.2,\quad
0.2\times(0.08)^2 = 0.2\times6.4\times10^{-3} = 1.28\times10^{-3}
\]
\[
1.28\times10^{-3}\times4\pi\times10^6 = 5120\pi\;\text{(not }5120\pi\times10^1\text{)}
\]
\[
|V_\text{emf}| = \frac{\sqrt{6}}{4}\times5120\pi = 1280\pi\sqrt{6}
\]
\[
1280\times\pi\approx4021,\quad 4021\times\sqrt{6}\approx4021\times2.449\approx9847\;\text{V}
\]
\[
\boxed{|V_\text{emf}| = 1280\pi\sqrt{6}\approx9847\;\text{V}}
\]

\step{3} Apply Ohm's law \cs{p.1: EMF \& FARADAY'S LAW equivalent circuit}.
\[
I = \frac{V_\text{emf}}{R} = \frac{1280\pi\sqrt{6}}{1000} = 1.28\pi\sqrt{6}
\]
\[
\boxed{I = 1.28\pi\sqrt{6}\approx 9.85\;\text{A}}
\]

\textit{Compare: prior keys got $\approx1.57\,\text{A}$ by omitting the $2\pi$ conversion.
The factor of $\approx6.28$ difference is exactly $2\pi$.}

\vspace{8pt}

%──────────────────────────────────────────────────────────────────────
\colorbox{purplebg}{\textbf{\color{purple}\large Question 2 --- Moving Wire (12 pts)}}
%──────────────────────────────────────────────────────────────────────

\vspace{4pt}
Wire along $\hat{y}$, $a=0.15\,\text{m}$, $\mathbf{u}=0.12\hat{x}\,\text{m/s}$,
$\mathbf{B}=(0.1\hat{x}-0.05\hat{z})\,\text{T}$.

\cs{p.1: EMF \& FARADAY'S LAW (Motional EMF)}: $V_\text{emf}=\displaystyle\int(\mathbf{u}\times\mathbf{B})\cdot d\mathbf{l}$.

\step{1} Compute $\mathbf{u}\times\mathbf{B}$ \cs{p.1: CROSS PRODUCT (determinant)}.
\[
\mathbf{u}\times\mathbf{B}=
\begin{vmatrix}\hat{x}&\hat{y}&\hat{z}\\0.12&0&0\\0.1&0&-0.05\end{vmatrix}
=\hat{x}(0\cdot(-0.05)-0\cdot0)
 -\hat{y}(0.12\cdot(-0.05)-0\cdot0.1)
 +\hat{z}(0.12\cdot0-0\cdot0.1)
\]
\[
= \hat{x}(0) - \hat{y}(-0.006) + \hat{z}(0) = 0.006\hat{y}\;\text{V/m}
\]

\step{2} Integrate along wire ($d\mathbf{l}=\hat{y}\,dy$).
\[
V_\text{emf}=\int_0^{0.15}0.006\,dy = 0.006\times0.15
\]
\[
\boxed{V_\text{emf} = 9\times10^{-4}\;\text{V} = 0.9\;\text{mV}}
\]

\textit{Only the $-0.05\hat{z}$ component of $\mathbf{B}$ contributes;
the $\hat{x}\times\hat{x}=0$ term drops out.
Common error: forgetting to convert cm/s $\to$ m/s (factor-of-100 mistake).}

\vspace{8pt}

%──────────────────────────────────────────────────────────────────────
\colorbox{coralbg}{\textbf{\color{coral}\large Question 3 --- Conduction \& Displacement Current (12 pts)}}
%──────────────────────────────────────────────────────────────────────

\vspace{4pt}
$I_c=2\sin(\omega t)\,\mu\text{A}$, cross-section $A$, $\sigma=10^{-4}\,\text{S/m}$, $\varepsilon_r=26\pi$.

\medskip

\textbf{(a) Electric field $E$.} \hfill\textit{(6 pts)}

\cs{p.1: DISPLACEMENT \& CONDUCTION CURRENT}: $J_c=\sigma E$, $I_c=J_c\cdot A$.

\[
J_c = \frac{I_c}{A}, \qquad E = \frac{J_c}{\sigma} = \frac{I_c}{\sigma A}
= \frac{2\times10^{-6}\sin(\omega t)}{10^{-4}\cdot A}
\]
\[
\boxed{E = \frac{0.02}{A}\sin(\omega t)\;\text{V/m}}
\]
\textit{Area $A$ remains in the answer — the problem gives no numeric value for it.}

\medskip

\textbf{(b) Displacement current $I_d$ at $\omega=10^9\,\text{rad/s}$, $\omega t=\pi/3$.} \hfill\textit{(6 pts)}

\cs{p.1: DISPLACEMENT \& CONDUCTION CURRENT}: $J_d=\varepsilon\,\partial E/\partial t$, $I_d=J_d\cdot A$.\\
\cs{p.1: $\partial B/\partial t$ TABLE}: $\partial[\sin(\omega t)]/\partial t=\omega\cos(\omega t)$.\\
\cs{p.1: UNIT CIRCLE}: $\cos(\pi/3)=\tfrac{1}{2}$.

\step{1} Compute $\varepsilon=\varepsilon_r\varepsilon_0$.
\[
\varepsilon = 26\pi\cdot\frac{1}{36\pi}\times10^{-9}
= \frac{26}{36}\times10^{-9}
= \frac{13}{18}\times10^{-9}\;\text{F/m}
\]

\step{2} $I_d = \varepsilon A\cdot\dfrac{\partial E}{\partial t}
= \varepsilon A\cdot\dfrac{0.02}{A}\cdot\omega\cos(\omega t)$
(the $A$ cancels).
\[
I_d = \frac{13}{18}\times10^{-9}\times 0.02\times10^9\times\frac{1}{2}
= \frac{13}{18}\times0.02\times\frac{1}{2}
= \frac{13}{1800}
\]
\[
\boxed{I_d = \frac{13}{1800}\approx 7.22\;\text{mA}}
\]

\vspace{8pt}

%──────────────────────────────────────────────────────────────────────
\colorbox{coralbg}{\textbf{\color{coral}\large Question 4 --- Charge Relaxation (10 pts)}}
%──────────────────────────────────────────────────────────────────────

\vspace{4pt}
$\sigma=10^{-4}\,\text{S/m}$, $\varepsilon_r=26\pi$.

\medskip

\textbf{(a) Relaxation time constant.} \hfill\textit{(5 pts)}

\cs{p.1: CHARGE RELAXATION}: $\tau_r=\varepsilon_r\varepsilon_0/\sigma$.
\[
\tau_r = \frac{26\pi\cdot\tfrac{1}{36\pi}\times10^{-9}}{10^{-4}}
= \frac{\tfrac{13}{18}\times10^{-9}}{10^{-4}}
= \frac{13}{18}\times10^{-5}\;\text{s}
\]
\[
\boxed{\tau_r = \frac{13}{18}\times10^{-5}\approx 7.22\;\mu\text{s}}
\]

\medskip

\textbf{(b) Time to decay to $e^{-3}$ of initial value.} \hfill\textit{(5 pts)}

\cs{p.1: CHARGE RELAXATION}: $\rho_v(t)=\rho_{v0}\,e^{-t/\tau_r}$.
\[
e^{-t/\tau_r}=e^{-3}
\;\Rightarrow\;\frac{t}{\tau_r}=3
\;\Rightarrow\;t=3\tau_r
\]
\[
\boxed{t = 3\times7.22\;\mu\text{s}\approx 21.7\;\mu\text{s}}
\]

\vspace{8pt}

%──────────────────────────────────────────────────────────────────────
\colorbox{purplebg}{\textbf{\color{purple}\large Question 5 --- Lossless EM Wave (30 pts)}}
%──────────────────────────────────────────────────────────────────────

\vspace{4pt}
$\mathbf{E}(z,t)=4\cos\!\bigl(6\pi\times10^8\,t-kz-\pi/6\bigr)\hat{y}$ V/m.
$\varepsilon_r=6$, $\mu_r=1$.

\medskip

\textbf{(a) Find $\omega$, $k$, $\lambda$, $\eta$.} \hfill\textit{(10 pts)}

\cs{p.2: PLANE WAVE PARAMETERS (LOSSLESS)}: read $\omega$ from the cosine argument;
chain $\omega\to u_p\to k\to\lambda\to\eta$.

\[
\omega = 6\pi\times10^8\;\text{rad/s}
\]
\[
k = \frac{\omega\sqrt{\varepsilon_r}}{c}
= \frac{6\pi\times10^8\cdot\sqrt{6}}{3\times10^8}
= 2\pi\sqrt{6}
\qquad\boxed{k = 2\pi\sqrt{6}\approx15.4\;\text{rad/m}}
\]
\[
\lambda = \frac{2\pi}{k} = \frac{2\pi}{2\pi\sqrt{6}} = \frac{1}{\sqrt{6}}
\qquad\boxed{\lambda = \frac{1}{\sqrt{6}}\approx0.408\;\text{m}}
\]
\[
\eta = \frac{\eta_0}{\sqrt{\varepsilon_r}} = \frac{120\pi}{\sqrt{6}}
= 20\pi\sqrt{6}
\qquad\boxed{\eta = 20\pi\sqrt{6}\approx154\;\Omega}
\]

\medskip

\textbf{(b) Find $\tilde{\mathbf{E}}(z)$.} \hfill\textit{(5 pts)}

\cs{p.1: PHASOR METHOD}: $\mathbf{E}=\operatorname{Re}\{\tilde{\mathbf{E}}\,e^{j\omega t}\}$; strip $e^{j\omega t}$, keep the rest.\\
\cs{p.2: PHASOR $\leftrightarrow$ TIME DOMAIN}: $Ae^{j\phi}e^{-j\beta z}\leftrightarrow A\cos(\omega t-\beta z+\phi)$.
\[
\boxed{\tilde{\mathbf{E}}(z) = 4\,e^{-j(kz+\pi/6)}\,\hat{y}\;\text{V/m}}
\]

\medskip

\textbf{(c) Find $\tilde{\mathbf{H}}(z)$.} \hfill\textit{(5 pts)}

\cs{p.2: DIRECTION OF PROPAGATION}: argument contains $-kz$ $\Rightarrow$ $\hat{k}=+\hat{z}$.\\
\cs{p.2: $\hat{k}\times\hat{E}$ --- FIND $\hat{H}$ DIRECTION}: $\hat{k}=+\hat{z}$, $\hat{E}=\hat{y}$ $\Rightarrow$
$\hat{z}\times\hat{y}=-\hat{x}$.\\
\cs{p.2: LOSSLESS vs.\ LOSSY --- E \& H}: $\tilde{\mathbf{H}}=\tfrac{1}{\eta}(\hat{k}\times\tilde{\mathbf{E}})$.
\[
\tilde{\mathbf{H}} = \frac{1}{\eta}\,\hat{z}\times\bigl(4e^{-j(kz+\pi/6)}\hat{y}\bigr)
= \frac{4}{\eta}\,e^{-j(kz+\pi/6)}\,(\hat{z}\times\hat{y})
\]
With $\eta=20\pi\sqrt{6}$: $\tfrac{4}{\eta}=\tfrac{4}{20\pi\sqrt{6}}=\tfrac{1}{5\pi\sqrt{6}}$.
\[
\boxed{\tilde{\mathbf{H}}(z)
= -\frac{1}{5\pi\sqrt{6}}\,e^{-j(kz+\pi/6)}\,\hat{x}\;\text{A/m}
\approx -0.026\,e^{-j(kz+\pi/6)}\,\hat{x}\;\text{A/m}}
\]

\medskip

\textbf{(d) Find $\mathbf{H}(z,t)$.} \hfill\textit{(5 pts)}

\cs{p.2: PHASOR $\leftrightarrow$ TIME DOMAIN}: $\operatorname{Re}\{Ae^{j\phi}e^{j\omega t}e^{-j\beta z}\}
= A\cos(\omega t-\beta z+\phi)$.
\[
\boxed{\mathbf{H}(z,t)
= -\frac{1}{5\pi\sqrt{6}}\cos\!\left(6\pi\times10^8\,t - 2\pi\sqrt{6}\,z - \frac{\pi}{6}\right)\hat{x}
\;\text{A/m}}
\]

\medskip

\textbf{(e) Average power density.} \hfill\textit{(5 pts)}

\cs{p.2: AVERAGE POWER DENSITY (POYNTING) lossless shortcut}: $\mathbf{S}_\text{av}=\tfrac{|E_0|^2}{2\eta}\,\hat{k}$.\\
$|E_0|=4\,\text{V/m}$, $\hat{k}=+\hat{z}$, $\eta=20\pi\sqrt{6}$:
\[
\mathbf{S}_\text{av} = \frac{16}{2\times20\pi\sqrt{6}}\,\hat{z}
= \frac{16}{40\pi\sqrt{6}}\,\hat{z}
= \frac{2}{5\pi\sqrt{6}}\,\hat{z}
= \frac{\sqrt{6}}{15\pi}\,\hat{z}
\]
\[
\boxed{\mathbf{S}_\text{av} = \frac{\sqrt{6}}{15\pi}\,\hat{z}\approx 0.052\;\hat{z}\;\text{W/m}^2}
\]

\vspace{8pt}

%──────────────────────────────────────────────────────────────────────
\colorbox{coralbg}{\textbf{\color{coral}\large Question 6 --- EM Wave in Distilled Water (16 pts)}}
%──────────────────────────────────────────────────────────────────────

\vspace{4pt}
\textbf{Given:} $f=\tfrac{2}{\pi}\,\text{GHz}$, $\sigma=10^{-4}\,\text{S/m}$, $\varepsilon_r=26\pi$, $\mu_r=1$.

\begin{notebox}
\textbf{Compute $\omega$ first \cs{p.2: WAVE PARAMS strip}: $\omega=2\pi f$}
\[
\omega = 2\pi\times\frac{2}{\pi}\times10^9 = 4\times10^9\;\text{rad/s}
\]
\textbf{Also useful:}
$\varepsilon' = \varepsilon_r\varepsilon_0
= 26\pi\cdot\tfrac{1}{36\pi}\times10^{-9}
= \tfrac{13}{18}\times10^{-9}\,\text{F/m}$
\end{notebox}

\medskip

\textbf{(a) Find $\varepsilon''/\varepsilon'$.} \hfill\textit{(6 pts)}

\cs{p.2: LOSSY MEDIA --- CLASSIFICATION}: $\varepsilon''/\varepsilon'=\sigma/(\omega\varepsilon_r\varepsilon_0)$.
\[
\frac{\varepsilon''}{\varepsilon'}
= \frac{10^{-4}}{4\times10^9\times\tfrac{13}{18}\times10^{-9}}
= \frac{10^{-4}}{\tfrac{52}{18}}
= \frac{10^{-4}\times18}{52}
= \frac{18}{52}\times10^{-4}
= \frac{9}{26}\times10^{-4}
\]
\[
\boxed{\frac{\varepsilon''}{\varepsilon'} = \frac{9}{26}\times10^{-4}\approx 3.46\times10^{-5}}
\]

\medskip

\textbf{(b) What type of medium?} \hfill\textit{(5 pts)}

\cs{p.2: LOSSY MEDIA --- CLASSIFICATION table}: threshold for low-loss is $\varepsilon''/\varepsilon'<0.01$.
\[
\frac{\varepsilon''}{\varepsilon'}\approx3.46\times10^{-5}\;\ll\;0.01
\]
\[
\boxed{\text{Low-loss dielectric.}}
\]
Justification: the ratio is far below the $10^{-2}$ threshold.

\medskip

\textbf{(c) Skin depth.} \hfill\textit{(5 pts)}

\cs{p.2: LOW-LOSS MEDIUM}: $\alpha\approx\tfrac{\sigma}{2}\sqrt{\mu/\varepsilon'}=\tfrac{\sigma\eta}{2}$.\\
\cs{p.2: SKIN DEPTH}: $\delta_s=1/\alpha$.

\step{1} Find $\eta$ \cs{p.2: WAVE PARAMS strip}.
\[
\eta = \frac{\eta_0}{\sqrt{\varepsilon_r}}
= \frac{120\pi}{\sqrt{26\pi}}
\approx\frac{376.99}{9.04}
\approx 41.7\;\Omega
\]

\step{2} Attenuation constant (low-loss approximation).
\[
\alpha \approx \frac{\sigma\eta}{2}
= \frac{10^{-4}\times41.7}{2}
\approx 2.09\times10^{-3}\;\text{Np/m}
\]

\step{3} Skin depth.
\[
\delta_s = \frac{1}{\alpha} = \frac{1}{2.09\times10^{-3}}
\]
\[
\boxed{\delta_s \approx 479\;\text{m}}
\]
\textit{$\approx$\,480\,m confirms this is essentially lossless at $\tfrac{2}{\pi}$\,GHz with $\sigma=10^{-4}$\,S/m.}

\vspace{12pt}
\hrule
\vspace{6pt}
\begin{center}
{\small\textbf{Score (if perfect):}
Q1: 20/20\quad Q2: 12/12\quad Q3: 12/12\quad Q4: 10/10\quad Q5: 30/30\quad Q6: 16/16
$=$ \textbf{100/100}}
\end{center}

\end{document}
